% Source: tex/chapters/ch05/06_テスト技法.tex
\subsubsection{技法の選択と組合せ}
実務では、一つの技法だけで十分な品質根拠を得ることは少ない。仕様ベースと構造ベースを組み合わせれば、外部振る舞いと内部網羅を両方確認できる。リスクが高い部分には手厚い技法を使い、低リスク部分には軽量な技法を使うなど、目的、費用、リスク、組織制約に応じて選ぶ。

\paragraph{機能的技法と構造的技法の組合せ}
シナリオベーステストは機能的観点を強く持つ。構造ベーステストは内部構造の観点を強く持つ。両者は対立ではなく補完関係である。仕様に書かれた動作を確認しつつ、実装上の分岐やデータ流れも確認することで、見落としを減らせる。

\paragraph{決定的テストとランダムテスト}
決定的テストは、基準に従って選ばれたテストケースを実行する。ランダムテストは、確率分布に従って入力を選ぶ。信頼性テストではランダム性が役立つことがあるが、境界や特定リスクを狙う場合には決定的な選択が有効である。
